\section{Open Questions}

The following are five genuinely open questions raised by --- but not answered
by --- this replication. Each includes an evidence basis and concrete,
free-endpoint-only next steps.

\subsection*{Q1: Extension to FLASH and proton-minibeam regimes}
\textbf{Basis.} The paper's dose-rate machinery (Sec 2.3, Fig 6) was
validated against Lehmann/Newman conventional dose rates. FLASH ($>40$~Gy/s)
and minibeam geometries lie orders of magnitude outside the calibration
window, yet the mechanistic MC backbone (fast/slow repair with rate
constants $\approx 2.1~\mathrm{h}^{-1}$/$0.26~\mathrm{h}^{-1}$) should
extrapolate cleanly if track structure is honestly resampled at those
instantaneous dose rates.

\textbf{Next steps.}
(1)~Modify \texttt{damageModel.basicXandIon} to accept an
\texttt{instantaneousDoseRate} override, seeding SDD generation at
$10^2$--$10^3$~Gy/s pulse profiles. (2)~Compare Medras-MC predicted misrepair
vs.\ published FLASH clonogenic data (Vozenin 2019 zebrafish/mouse
survival). (3)~Report divergence as a function of pulse duty cycle. No
external paid data needed; FLASH survival tables are in supplementary of
most FLASH papers.

\subsection*{Q2: Modern multi-pathway DNA-repair kinetics}
\textbf{Basis.} The shipped code hard-codes a two-component repair model
with rate constants $\lambda_f$ and $\lambda_s$ that lump NHEJ into
``fast'' and HR-plus-NHEJ-on-complex into ``slow''. Recent single-molecule
imaging (Britton 2020, Setiaputra 2019) resolves at least four
kinetically-distinct pathways with cell-cycle-phase and chromatin-state
dependence.

\textbf{Next steps.}
(1)~Fork \texttt{medrasrepair.py} and expose the rate-constant table as a
JSON config. (2)~Insert a 4-pathway model with published rate constants
from L\"offler 2018. (3)~Re-run the 23-condition sweep and compare
bi-exponential fit residuals. (4)~Check whether high-LET carbon predictions
shift enough to change the RBE-peak location.

\subsection*{Q3: LET-spectrum sensitivity (mono-energetic vs.\ realistic SOBP)}
\textbf{Basis.} This replication used mono-energetic proton tracks
1.77--29.78~keV/$\mu$m. Clinical spread-out-Bragg-peak (SOBP) beams
sample a wide LET distribution at each depth, with an especially sharp LET
rise in the distal fall-off region where RBE uncertainty matters most.

\textbf{Next steps.}
(1)~Generate synthetic SOBP LET spectra using TOPAS or the analytic
Bortfeld model. (2)~Run \texttt{damageModel} with per-voxel LET-weighted
averages vs.\ properly-sampled spectra. (3)~Quantify how much the
misrepair-vs-depth curve differs between mono-energetic and spectrum-realistic
inputs. Free code: TOPAS is open-source.

\subsection*{Q4: Cell-line-specific prediction outliers}
\textbf{Basis.} The 2021 paper's Fig~4 MID scatter shows systematic
outliers, but the paper does not diagnose them. Cell lines with abnormal
repair phenotypes (BRCA1/2 loss, DNA-PKcs mutation, ATM-null) should map to
specific Medras parameters (fast/slow branching ratio, complex-lesion
misrepair probability).

\textbf{Next steps.}
(1)~Curate a small panel of cell-survival curves from cell lines with
sequenced DDR-pathway variants (via cBioPortal or DepMap, both free).
(2)~Fit Medras parameters per line. (3)~Report which parameters correlate
with which mutation classes. (4)~Test whether the residuals shrink when a
mutation-class-specific parameter override is applied.

\subsection*{Q5: ML-driven parameter identification from clonogenic assays}
\textbf{Basis.} Currently Medras uses a single global parameter set.
Clonogenic assays measure survival at 5--10 doses per cell line; a modern
amortized-inference network (e.g., BayesFlow, sbi) could invert those
5--10 measurements to a full Medras parameter posterior per line in seconds.

\textbf{Next steps.}
(1)~Build a simulator wrapper that emits survival curves from Medras
parameter tuples. (2)~Train a neural posterior estimator on $10^4$
simulated parameter/curve pairs. (3)~Validate on held-out published
clonogenic assays. (4)~Compare inferred parameters across radioresistant
vs.\ radiosensitive lines. All-free stack: PyTorch + sbi + published
clonogenic tables from open supplementary materials.
